%%% Richard James Howe 2/2
%% Final Year Project poster presentation.
% This uses a package called 'baposter.cls' which can be
% found here: http://www.brian-amberg.de/uni/poster/
% This package has been written by Brian Amberg (baposter@brian-amberg.de)

\documentclass[a1paper,portrait]{baposter}

\usepackage{relsize}    % For \smaller
\usepackage{url}      % For \url
\usepackage{epstopdf}  % Included EPS files automatically converted to PDF to include with pdflatex

%%% Global Settings %%%%%%%%%%%%%%%%%%%%%%%%%%%%%%%%%%%%%%%%%%%%%%%%%%%%%%%%%%%

\graphicspath{{pix/}}  % Root directory of the pictures 
\tracingstats=2      % Enabled LaTeX logging with conditionals

%%% Color Definitions %%%%%%%%%%%%%%%%%%%%%%%%%%%%%%%%%%%%%%%%%%%%%%%%%%%%%%%%%

\definecolor{bordercol}{RGB}{40,40,40}
\definecolor{headercol1}{RGB}{186,215,230}
\definecolor{headercol2}{RGB}{80,80,80}
\definecolor{headerfontcol}{RGB}{0,0,0}
\definecolor{boxcolor}{RGB}{186,215,230}

%%%%%%%%%%%%%%%%%%%%%%%%%%%%%%%%%%%%%%%%%%%%%%%%%%%%%%%%%%%%%%%%%%%%%%%%%%%%%%%%
%%% Utility functions %%%%%%%%%%%%%%%%%%%%%%%%%%%%%%%%%%%%%%%%%%%%%%%%%%%%%%%%%%

%%% Save space in lists. Use this after the opening of the list %%%%%%%%%%%%%%%%
\newcommand{\compresslist}{
  \setlength{\itemsep}{1pt}
  \setlength{\parskip}{0pt}
  \setlength{\parsep}{0pt}
}

%%%%%%%%%%%%%%%%%%%%%%%%%%%%%%%%%%%%%%%%%%%%%%%%%%%%%%%%%%%%%%%%%%%%%%%%%%%%%%%
%%% Document Start %%%%%%%%%%%%%%%%%%%%%%%%%%%%%%%%%%%%%%%%%%%%%%%%%%%%%%%%%%%%
%%%%%%%%%%%%%%%%%%%%%%%%%%%%%%%%%%%%%%%%%%%%%%%%%%%%%%%%%%%%%%%%%%%%%%%%%%%%%%%

\begin{document}
\typeout{Poster rendering started}

%%% Setting Background Image %%%%%%%%%%%%%%%%%%%%%%%%%%%%%%%%%%%%%%%%%%%%%%%%%%
\background{
  \begin{tikzpicture}[remember picture,overlay]%
  \draw (current page.north west)+(-2em,2em) node[anchor=north west]
  {\includegraphics[height=1.1\textheight]{background}};
  \end{tikzpicture}
}

\begin{poster}{
  grid=false,
  borderColor=bordercol,
  headerColorOne=headercol1,
  headerColorTwo=headercol2,
  headerFontColor=headerfontcol,
  boxColorOne=boxcolor,
  headershape=roundedright,
  headerfont=\Large\sf\bf,
  textborder=rectangle,
  background=user,
  headerborder=open,
  boxshade=plain
}
{
  Eye Catcher, empty if option eyecatcher=false - unused
}
{\sf\bf
  A VHDL based educational computer (2/2).
}
{
  \vspace{1em} Richard James Howe : 082060589\\
  {\smaller howe.r.j.89@gmail.com}
}
{
\setlength\fboxsep{0pt}
\setlength\fboxrule{0.5pt}
  \fbox{
    \begin{minipage}{10em}
      \includegraphics[width=10em,height=4em]{aston.jpg}
    \end{minipage}
  }
}

\headerbox{Abstract}{name=abstract,column=0,row=0}{
\smaller                          % Make the whole text smaller
}


\headerbox{Github}{name=github,column=0,below=abstract}{
\smaller                          % Make the whole text smaller
}

\headerbox{Future plans}{name=futureplans,column=0,below=github, above=bottom}{
\smaller                          % Make the whole text smaller
} 

\headerbox{How it works}{name=howitworks,span=2,column=1,row=0}{
\smaller

The diagram bellow shows how data moves through the H2 (The name of this
CPU), the processor can execute one instruction per clock cycle and has room for
further customization.

\begin{center}
    \includegraphics[scale=0.12]{h2.png}
\end{center}
}

\headerbox{Extensibility}
{name=firmware,span=2,column=1,below=howitworks,above=bottom}{
\smaller

Both the CPU core and the system itself are extensible, you can modify both. The CPU by adding
more instructions, the system as a whole by adding new devices to the interface list. I have
added VGA support by including a VGA core available on \url{opencore.org}.

}

\end{poster}
\end{document}
